\section{Global convergence: mixed benchmarks, inward models, and homogeneous geometry}\label{art:global}

We now turn the comparison-energy framework of Section~\ref{art:energy} into
concrete global convergence theorems on $\R^n$. To make the mechanism
completely transparent before introducing general sublevel-band definitions, we
begin by resolving the open benchmark $F(u,v)=u^2+|v|$ of Guo and
Lewis~\cite{GL18} directly in a few lines, and then extract the general
one-point inward principle that powers our $n$-dimensional theorems.

\subsection{Warm-up: whole-sequence convergence on $F(u,v)=u^2+|v|$}

On $F(u,v)=u^2+|v|$, Guo and Lewis~\cite[Corollary~4.4]{GL18} proved via smooth
surrogates that $\liminf_{k\to\infty}|v_k|=0$, leaving open whether the full
sequence $(u_k,v_k)$ converges to $(0,0)$. With Lemma~\ref{art:reference-energy}
and Theorem~\ref{art:reciprocal-growth} in hand, whole-sequence convergence
follows directly from a two-line radial secant estimate.

\begin{corollary}[Whole-sequence convergence on $u^2+|v|$]\label{art:mixed-benchmark}
For $F(u,v)=u^2+|v|$, every infinite admissible full-BFGS run, with arbitrary
$H_0\succ0$ and any fixed $0<c_1<c_2<1$, satisfies
\[
 F(u_k,v_k)\longrightarrow0,\qquad (u_k,v_k)\longrightarrow(0,0).
\]
\end{corollary}
\begin{proof}
Set $b=F(u_0,v_0)>0$; Armijo keeps $x_k=(u_k,v_k)$ inside the compact sublevel
$K=\{F\le b\}$, where gradients $g_k=(2u_k,\operatorname{sign}v_k)$ satisfy
$\norm{g_k}^2=4u_k^2+1\le4b+1$ and $\sum_kd_k<\infty$. Consider the reference
energy $E_k=x_k^TB_kx_k$ centered at $x_*=(0,0)$. For any two differentiable
points $x=(u,v)$ and $z=(r,w)$ with $y=\nabla F(z)-\nabla F(x)$ and
$t=y^T(z-x)>0$:
\begin{itemize}
\item if $v$ and $w$ have the same sign, then $y=(2(r-u),0)$, so
$t=2(r-u)^2$ and $(y^Tz)^2=4(r-u)^2r^2=2r^2t\le2F(z)t$;
\item if $v$ and $w$ have opposite signs, then with $S=|v|+|w|>0$ we have
$t=2(r-u)^2+2S$ and $y^Tz=2(r-u)r+2|w|$, so Cauchy--Schwarz gives
\[
 (y^Tz)^2\le\bigl(2(r-u)^2+2S\bigr)\left(2r^2+\frac{2|w|^2}{S}\right)
 \le2(r^2+|w|)t=2F(z)t.
\]
\end{itemize}
In either case, $R_k=(y_k^Tx_{k+1})^2/t_k\le2F(x_{k+1})\le2b$ for every step
$k$. By~\eqref{art:energy-budget} of Lemma~\ref{art:reference-energy}, a
positive limiting value $\lim_kF(x_k)=\eta>0$ would imply
$\eta^2\sum_{k<N}q_k^{-1}\le E_0+2bN$, contradicting
Theorem~\ref{art:reciprocal-growth}. Hence $u_k^2+|v_k|\to0$ and
$(u_k,v_k)\to(0,0)$.\qed
\end{proof}

Why did the cross-kink jump $y_k=(2(u_{k+1}-u_k),\pm2)$ satisfy
$(y_k^Tx_{k+1})^2\le2bt_k$ even when $s_k\to0$ makes $\norm{y_k}^2/t_k\to+\infty$?
Because along every inward segment toward the origin, $|v|$ scales exactly
linearly while $u^2$ contributes only a quadratic Taylor error~\eqref{art:mixed-inward-remainder}
(an explicit admissible run with $H_k=\operatorname{diag}(1,2^{-k})\to\operatorname{diag}(1,0)$
is exhibited in Proposition~\ref{art:examples-mixed}). We now show that this
one-point inward property controls two-point radial secants for arbitrary finite
convex functions on $\R^n$.

\subsection{From one-point inward models to two-point secant control}

Fix a minimizer $x_*$ of finite convex $F:\R^n\to\R$, write $F_*=F(x_*)$, and
let $K=\{x:F(x)\le b\}$ be compact with $b>F_*$. For differentiable points
$x,z\in K$, write
\[
 \begin{gathered}
 r_x=x-x_*,\qquad g=\nabla F(x),\qquad h=\nabla F(z),\\
 T=(h-g)^T(z-x),\qquad N=(h-g)^Tr_z,
 \end{gathered}
\]
so $T$ is the secant curvature $t$ of Section~\ref{art:energy} and $N$ is the
radial projection of the gradient jump at $z$. The \emph{radial secant
condition} requires control of $N^2/T$ across all descending differentiable
pairs on each positive sublevel band:
\begin{equation}\label{art:radial-condition}
 \begin{gathered}
 \text{for every }\eta>0\text{ there exists }C_\eta\ge0\text{ such that}\\
 N^2\le C_\eta T
 \quad\text{whenever }F_*+\eta\le F(z)\le F(x)\text{ and }T>0.
 \end{gathered}
\end{equation}
Rather than checking pairs of points $(x,z)$ across different smooth regions by
hand, the \emph{inward model condition} tests $F$ at a single differentiable
point $x$ along the inward ray toward $x_*$:
\begin{equation}\label{art:inward-condition}
 \begin{gathered}
 \text{for every }\eta>0\text{ there exist }M_\eta\ge0\text{ and }
       \theta_\eta\in(0,1/2]\text{ such that}\\
 0\le F(x-\tau r_x)-F(x)+\tau\nabla F(x)^Tr_x
       \le\frac{M_\eta}{2}\tau^2\quad(0\le\tau\le\theta_\eta)
 \end{gathered}
\end{equation}
at every differentiable $x\in K$ with $F(x)\ge F_*+\eta$. By convexity, the
inward sample $x-\tau r_x$ remains in $K$ and need not be differentiable; $\tau$
is a dimensionless fraction of the ray to $x_*$, and both $M_\eta$ and
$\theta_\eta$ may depend on the band $\eta>0$.

\begin{theorem}[Inward models control radial secants]\label{art:radial-equivalence}
For any finite convex $F:\R^n\to\R$ and compact sublevel $K=\{F\le b\}$, the
one-point inward condition~\eqref{art:inward-condition} implies the two-point
radial secant condition~\eqref{art:radial-condition}.
\end{theorem}
\begin{proof}
Choose $\Lambda,R>0$ bounding all subgradient norms and $\norm{x-x_*}$ on $K$
(existence of $\Lambda$ was established in the proof of
Theorem~\ref{art:certificate}). Define the Bregman supporting-plane
residual~\cite{BB}
\[
 D_F(z,x)=F(z)-F(x)-\nabla F(x)^T(z-x)\ge0,
\]
which requires differentiability only at its second argument and satisfies
$T=D_F(z,x)+D_F(x,z)$ when both $x$ and $z$ are differentiable.

Fix $\eta>0$ and let $M=M_\eta$, $\theta=\theta_\eta$ be the constants
from~\eqref{art:inward-condition}. Let $x,z\in K$ be differentiable with
$F_*+\eta\le F(z)\le F(x)$ and $T>0$. If $N\ge0$, comparing the supporting
inequality $F(z-\tau r_z)\ge F(x)+\nabla F(x)^T(z-\tau r_z-x)$ at $x$ with the
upper inward bound~\eqref{art:inward-condition} at $z$ yields
\[
 \tau N\le D_F(z,x)+\tfrac12M\tau^2\qquad(0\le\tau\le\theta).
\]
If $N<0$, reversing the roles of $x$ and $z$ (applying the supporting plane at
$z$ and the inward bound at $x$, where $F(x)\ge F(z)\ge F_*+\eta$) gives
$(g-h)^Tr_x=T-N\ge|N|$ and
\[
 \tau(T-N)\le D_F(x,z)+\tfrac12M\tau^2\qquad(0\le\tau\le\theta).
\]
In both cases there exist $A\in[|N|,2\Lambda R]$ and $D\in\{D_F(z,x),D_F(x,z)\}$
satisfying $\tau A\le D+\tfrac12M\tau^2$ for all $\tau\in[0,\theta]$. Optimizing
over $\tau$: if $M>0$ and $A\le M\theta$, setting $\tau=A/M$ gives
$A^2\le2MD$; if $A>M\theta$ (including the case $M=0$), setting $\tau=\theta$
gives $D\ge\tau A-\tfrac12M\tau^2\ge\tfrac12\theta A$ and hence
$A^2\le(4\Lambda R/\theta)D$. Since $D\le T$, we conclude that
\begin{equation}\label{art:inward-to-radial}
 N^2\le \max\{2M,4\Lambda R/\theta\}\,T.\qquad\qed
\end{equation}
\end{proof}

Chapter~4 of the companion proves that the converse implication also holds on
compact sublevels, so~\eqref{art:radial-condition} and~\eqref{art:inward-condition}
are in fact equivalent characterizations of radial comparison geometry.

\begin{corollary}[Radial convergence]\label{art:radial-convergence}
Under either~\eqref{art:radial-condition} or~\eqref{art:inward-condition} on a
compact sublevel $K=\{F\le b\}$, every admissible run starting at $x_0\in K$
satisfies $F(x_k)\to F_*$ and $\dist(x_k,\argmin F)\to0$, with $x_k\to x_*$
when $x_*$ is the unique minimizer.
\end{corollary}
\begin{proof}
By Theorem~\ref{art:radial-equivalence},~\eqref{art:radial-condition} holds on
$K$. Take the rank-one field $W_\eta(x)=r_xr_x^T$ on $K_\eta$. Convexity gives
$g^Tr_x\ge F(x)-F_*\ge\eta$, so~\eqref{art:visibility} holds with $m_\eta=\eta^2$.
Every eligible pair $(x,z)$ in $K_\eta$ has $F_*+\eta\le F(z)<F(x)$ and $T=t>0$,
so~\eqref{art:radial-condition} gives~\eqref{art:secant-control}. Finally,
$r_z=r_x+s$ gives $\Pi_sr_z=\Pi_sr_x$, so~\eqref{art:transport} holds with
$L_\eta=0$. Theorem~\ref{art:certificate} gives the result.\qed
\end{proof}

Geometrically,~\eqref{art:inward-condition} demands quadratic first-order
accuracy \emph{only along rays toward $x_*$}: in smooth regions its
infinitesimal form is $r_x^T\nabla^2F(x)r_x$ rather than $\lambda_{\max}(\nabla^2F(x))$,
while its finite-difference form also rules out uncompensated radial derivative
jumps. Transverse kinks whose switching surfaces are cones aimed at $x_*$ have
zero inward error regardless of how wildly their gradients jump in angular
directions.

\subsection{Centered smooth-plus-homogeneous objectives in $\R^n$}

Any finite convex positively homogeneous function $\gamma:\R^n\to\R$ of degree
one satisfies $\gamma((1-\tau)r)=(1-\tau)\gamma(r)$ for $0\le\tau\le1$, so its
inward Taylor remainder at $x_*$ vanishes identically regardless of angular
nonsmoothness. Adding an arbitrary locally $C^{1,1}$ convex function $\psi$
contributes only a smooth quadratic remainder, even when neither $\psi$ nor
$\gamma(\cdot-x_*)$ attains a minimum at $x_*$ on its own.

\begin{theorem}[Centered smooth-plus-homogeneous objectives]\label{art:centered}
Let $F:\R^n\to\R$ be finite and convex, and let $x_*$ minimize $F$. Let
$U\subseteq\R^n$ be an open convex set, let $\psi:U\to\R$ be convex and locally
$C^{1,1}$, and let $\gamma:\R^n\to\R$ be finite, convex, and positively
homogeneous of degree one. Suppose
\[
 F(x)=\psi(x)+\gamma(x-x_*)\qquad(x\in U).
\]
Then every admissible run whose initial sublevel $K=\{x:F(x)\le F(x_0)\}$ is
compact and contained in $U$ satisfies $F(x_k)\to F(x_*)$ and
$\dist(x_k,\argmin F)\to0$; if $x_*$ is the unique minimizer, then $x_k\to x_*$.
\end{theorem}
\begin{proof}
Since $g_0\ne0$, convexity gives $F(x_0)>F(x_*)$. Because $K\subset U$ is
compact and convex, $\nabla\psi$ has a uniform Lipschitz constant $L$ on $K$:
covering $K$ by finitely many open sets on which $\nabla\psi$ is Lipschitz and
subdividing any segment in $K$ into pieces shorter than a Lebesgue number of
that cover bounds the variation of $\nabla\psi$ along the segment by the
maximum local Lipschitz constant. Integrating along the inward segment from $x$
to $x-\tau r_x$ therefore gives
\[
 \left|\psi(x-\tau r_x)-\psi(x)+\tau\nabla\psi(x)^Tr_x\right|
 \le\tfrac12L\tau^2\norm{r_x}^2\qquad(0\le\tau\le1).
\]
At any differentiable point $x\in K$ of $F$, smoothness of $\psi$ implies that
$\gamma$ is differentiable at $r_x=x-x_*$, and Euler's identity for degree-one
homogeneous functions gives $\nabla\gamma(r_x)^Tr_x=\gamma(r_x)$, so
$\gamma((1-\tau)r_x)-\gamma(r_x)+\tau\nabla\gamma(r_x)^Tr_x=0$. Hence
\eqref{art:inward-condition} holds on $K$ with uniform constants
$M_\eta=L\max_{x\in K}\norm{r_x}^2$ and $\theta_\eta=1/2$ (the lower bound $0$
follows from convexity of $F$), and Corollary~\ref{art:radial-convergence}
completes the proof.\qed
\end{proof}

Notice that Theorem~\ref{art:centered} recovers
Corollary~\ref{art:mixed-benchmark} immediately in $\R^2$ by taking
$\psi(u,v)=u^2$, $\gamma(u,v)=|v|$, $x_*=(0,0)$, and $U=\R^2$, and extends it
to arbitrary smooth-plus-norm or smooth-plus-support-function sums in $\R^n$
such as $\psi(x)+\norm{Ax}_1$ or $\psi(x)+\max_ia_i^Tx$ whenever the composite
minimizer $x_*$ centers the homogeneous term (or locally after active-sign
splitting, as we exploit for the Lasso in Section~\ref{art:local}).

\subsection{Coercive positively homogeneous functions, norms, and linear maxima}

When $\psi=0$, the inward remainder vanishes everywhere on $\R^n$. Imposing
coercivity prevents rays of zero or negative slope and turns value convergence
into global iterate convergence from every starting point, with no smoothness or
finite-piece hypothesis.

\begin{theorem}[Coercive positively homogeneous objectives]\label{art:homogeneous}
Let $F:\R^n\to\R$ ($n\ge1$) be finite, convex, positively homogeneous of degree
one, and coercive. Then $0$ is the unique minimizer of $F$, and every infinite
admissible full-BFGS run from any differentiable $x_0\ne0$, with any
$H_0\succ0$ and fixed $0<c_1<c_2<1$, satisfies
\[
 F(x_k)\longrightarrow0,\qquad x_k\longrightarrow0.
\]
\end{theorem}
\begin{proof}
Positive homogeneity gives $F(0)=0$. If $F(x)\le0$ for some $x\ne0$, then
$F(tx)=tF(x)\le0$ for all $t>0$, contradicting coercivity as $t\to\infty$. Thus
$F(x)>0$ for all $x\ne0$, and continuity on the compact Euclidean unit sphere
gives
\[
 \mu=\min_{\norm{x}=1}F(x)>0,\qquad F(x)\ge\mu\norm{x}\quad(x\in\R^n).
\]
Every sublevel of $F$ is therefore compact and $0$ is the unique minimizer.
Applying Theorem~\ref{art:centered} with $x_*=0$, $U=\R^n$, $\psi=0$, and
$\gamma=F$ gives $F(x_k)\to0$, and $F(x_k)\ge\mu\norm{x_k}$ gives $x_k\to0$.\qed
\end{proof}

\begin{corollary}[Norms, gauges, and coercive linear maxima]\label{art:homogeneous-examples}
The global convergence conclusion of Theorem~\ref{art:homogeneous} holds on
$\R^n$ for:
\begin{enumerate}
\item every norm on $\R^n$ (in particular $\norm{x}_1$ and $\norm{x}_\infty$);
\item every finite convex degree-one gauge whose unit sublevel $\{x:F(x)\le1\}$
is bounded; and
\item every finite linear maximum $F(x)=\max_{1\le i\le m}a_i^Tx$ whose zero
sublevel $\{x:F(x)\le0\}$ is bounded.
\end{enumerate}
\end{corollary}
\begin{proof}
Every norm is finite, convex, degree-one homogeneous, and positive on the
Euclidean unit sphere, hence satisfies $F(x)\ge\mu\norm{x}$ with $\mu>0$. For a
degree-one gauge with $\{F\le1\}$ bounded or a finite linear maximum with
$\{F\le0\}$ bounded, no nonzero $x$ can satisfy $F(x)\le0$, since the entire ray
$\{tx:t\ge0\}$ would then lie inside that bounded sublevel; compactness of the
unit sphere again yields $F(x)\ge\mu\norm{x}$ with $\mu>0$, so $F$ is coercive
and Theorem~\ref{art:homogeneous} applies.\qed
\end{proof}

Corollary~\ref{art:homogeneous-examples} goes well beyond isolated kinks: in
dimension $n\ge2$, $\norm{x}_1=\sum_i|x_i|$ is nondifferentiable along every
coordinate hyperplane, and $\norm{x}_\infty=\max_i|x_i|$ has gradient jumps
along every diagonal hyperplane $|x_i|=|x_j|$ where two largest coordinates tie,
recovering and vastly generalizing the Euclidean-norm theorem of Guo and
Lewis~\cite[Corollary~4.2]{GL18} and the one-dimensional $|x|$ analysis of Lewis
and Overton~\cite[Section~5.1]{LO} without any symmetry requirement. Moreover,
Section~\ref{art:failure} (Theorem~\ref{art:failure_polyhedral}) shows that
absence of nonzero affine intercepts in $\max_ia_i^Tx$ is sharp: adding a single
nonzero offset in $\R^2$ produces a coercive four-piece polyhedral objective on
which full BFGS converges to a nonstationary corner. (When coercivity itself is
dropped on a positively homogeneous function, Chapter~7 of the companion proves
the global value dichotomy $F(x_k)\to0$ or $F(x_k)\to-\infty$.)
